\documentclass[10pt]{article}

% ---------- Document Identity (update these only) ----------
\newcommand{\DocStage}{}
\newcommand{\DocTitle}{Your Road to Tier One SOF}
\newcommand{\ProgramName}{182-Week Civilian-to-Selection Readiness Pipeline}
\newcommand{\ProgramSubtitle}{}

% ---------- Page ----------
\usepackage[legalpaper,left=0.5in,right=0.5in,top=0.6in,bottom=0.45in]{geometry}

% ---------- Typography ----------
\usepackage[T1]{fontenc}
\usepackage[utf8]{inputenc}
\usepackage{libertinus}
\usepackage{microtype}
\usepackage{parskip}
\usepackage{ragged2e}
\usepackage{textcomp}
\AtBeginDocument{\color{black}}
\AtBeginDocument{\pagecolor{white}}
\AtBeginDocument{\justifying}

% ---------- Landscape pages ----------
\usepackage{pdflscape}

% ---------- Color ----------
\usepackage[table]{xcolor}
\definecolor{StageBlue}{HTML}{1F4E79}
\definecolor{StageBlueDark}{HTML}{163A5A}
\definecolor{LightGray}{HTML}{F2F4F7}
\definecolor{MidGray}{HTML}{D9DEE5}
\definecolor{RuleGray}{HTML}{C7CED8}
\definecolor{HeaderInk}{HTML}{000000}
\definecolor{Ink}{HTML}{000000}

% ---------- Utility ----------
\usepackage{enumitem}
\sloppy
\setlength{\emergencystretch}{2em}

% ---------- Boxes / UI ----------
\usepackage[most]{tcolorbox}

\tcbset{
  charterTitle/.style={
    enhanced,
    colback=StageBlue!4,
    colframe=StageBlue,
    boxrule=1.25pt,
    arc=3pt,
    left=16pt,right=16pt,top=14pt,bottom=14pt,
    borderline north={3.2pt}{0pt}{StageBlue},
    borderline west={4pt}{0pt}{StageBlue},
    borderline south={1.25pt}{0pt}{StageBlue},
    drop shadow={black!25!white},
  },
  charterRule/.style={
    enhanced,
    colback=LightGray,
    colframe=RuleGray,
    boxrule=0.85pt,
    arc=3pt,
    left=12pt,right=12pt,top=10pt,bottom=10pt,
    borderline west={3pt}{0pt}{StageBlue},
  },
  charterBadge/.style={
    enhanced,
    colback=StageBlue,
    coltext=white,
    colframe=StageBlueDark,
    boxrule=0.6pt,
    arc=2pt,
    left=6pt,right=6pt,top=3pt,bottom=3pt,
  }
}
\newtcbox{\Badge}{nobeforeafter,charterBadge}

% ---------- Lists ----------
\setlist[itemize]{leftmargin=1.5em, itemsep=0.25em, topsep=0.25em}
\setlist[enumerate]{leftmargin=1.8em, itemsep=0.25em, topsep=0.25em}

% ---------- TOC ----------
\usepackage{tocloft}
\setcounter{tocdepth}{3}
\setlength{\cftbeforesecskip}{3pt}
\setlength{\cftbeforesubsecskip}{2pt}
\setlength{\cftbeforesubsubsecskip}{0pt}
\renewcommand{\cftsecfont}{\bfseries}
\renewcommand{\cftsecpagefont}{\bfseries}
\renewcommand{\cftsubsecfont}{\normalfont}
\renewcommand{\cftsubsecpagefont}{\normalfont}
\renewcommand{\cftsubsubsecfont}{\normalfont}
\renewcommand{\cftsubsubsecpagefont}{\normalfont}
\setlength{\cftsecindent}{0pt}
\setlength{\cftsubsecindent}{1.25em}
\setlength{\cftsubsubsecindent}{2.8em}
\setlength{\cftsecnumwidth}{2.1em}
\setlength{\cftsubsecnumwidth}{2.8em}
\setlength{\cftsubsubsecnumwidth}{3.6em}

% Remove the built-in TOC heading so only the custom heading is shown
\makeatletter
\renewcommand{\tableofcontents}{\@starttoc{toc}}
\makeatother

% ---------- Header/Footer: page number only ----------
\usepackage{fancyhdr}
\pagestyle{fancy}
\fancyhf{}
\renewcommand{\headrulewidth}{0pt}
\renewcommand{\footrulewidth}{0pt}
\fancyfoot[C]{\thepage}

% ---------- Section formatting ----------
\usepackage{titlesec}

\titleformat{\section}
  {\Large\bfseries\color{Ink}}
  {\thesection}{0.6em}{}
  [\vspace{0.25em}\textcolor{StageBlue}{\rule{\linewidth}{0.8pt}}]

\titleformat{\subsection}
  {\large\bfseries\color{Ink}}
  {\thesubsection}{0.6em}{}

\titleformat{\subsubsection}
  {\normalsize\bfseries\color{Ink}}
  {\thesubsubsection}{0.6em}{}

\titlespacing*{\section}{0pt}{0.95em}{0.35em}
\titlespacing*{\subsection}{0pt}{0.65em}{0.25em}
\titlespacing*{\subsubsection}{0pt}{0.55em}{0.2em}

% ---------- Table engine ----------
\usepackage{tabularray}
\UseTblrLibrary{booktabs}

% ---------- tabularray: GLOBAL table treatment ----------
\SetTblrInner{
  cells = {fg=black, bg=white, valign=t},
}

% ---------- Header helpers ----------
\newcommand{\Hdr}[1]{\shortstack[c]{\strut #1\strut}}

% ---------- Lines ----------
\newcommand{\TitleLine}{\textcolor{StageBlue}{\rule{0.98\linewidth}{0.95pt}}}
\newcommand{\SubTitleLine}{\textcolor{RuleGray}{\rule{0.98\linewidth}{0.65pt}}}

% ---------- Continued on next page ----------
\newcommand{\ContinuedOnNextPageInline}{%
  \par
  \vfill
  \noindent\textcolor{RuleGray}{\rule{\linewidth}{1pt}}\vspace{-2pt}
  \noindent\hfill\textit{Continued on next page}\par\vspace{-10pt}
  \noindent\textcolor{RuleGray}{\rule{\linewidth}{1pt}}
  \clearpage
}

% ---------- PDF hyperlinks ----------
\usepackage[hidelinks]{hyperref}
\hypersetup{
  colorlinks=false,
  pdfborder={0 0 0},
  linktoc=all,
  pdftitle={\DocTitle},
  pdfauthor={\ProgramName}
}

% =========================================================
\begin{document}

\section*{10.A — Macro Ladder Phases 00–13}
\phantomsection
\addcontentsline{toc}{section}{10.A — Macro Ladder Phases 00–13}

\subsection*{10.A.1 Macro Ladder Doctrine Program-Wide Rules}
\phantomsection
\addcontentsline{toc}{subsection}{10.A.1 Macro Ladder Doctrine Program-Wide Rules}

\begin{itemize}
\item Calorie posture progression across the continuum
  \begin{itemize}
  \item Phase 00–Phase 01: default is a safe, adherence-forward deficit that prioritizes habit formation and consistent execution under start-from-zero constraints.
  \item Phase 02–Phase 06: default is a moderated deficit to near-maintenance posture to protect strength development, running tolerance, and recovery so performance evidence remains valid and repeatable.
  \item Phase 07–Phase 13: default is near-maintenance to slight surplus posture with strict guardrails to protect high-specificity work, high-consequence testing, and validity in protected windows.
  \end{itemize}

\item Protein doctrine hybrid anchor expressed as grams per day ranges
  \begin{itemize}
  \item Anchor A basis: target bodyweight $\approx$ 180 lb.
  \item Anchor B basis: estimated lean body mass (LBM) derived from start-state body fat estimate (uncertain).
  \item Hybrid basis rule: use the higher of Anchor A basis and Anchor B basis as the protein basis; apply a conservative ceiling to avoid unnecessary complexity.
  \item Multiplier rule: 1.0–1.2 g/lb of the chosen basis.
  \item Program-wide protein target range (initial): 180–216 g/day, with bounded downward adjustment only as bodyweight decreases and intake feasibility demands, while preserving satiety, lean mass retention, and recovery support.
  \end{itemize}

\item Fat doctrine percent-based floor rule
  \begin{itemize}
  \item Fat floor \textbf{MUST} remain 20–30\% of total calories.
  \item Purpose: sustainability, satiety, adherence, and avoidance of overly aggressive restriction that degrades recovery and increases drift risk.
  \end{itemize}

\item Carbohydrate doctrine grams per day ranges and timing notes
  \begin{itemize}
  \item Carbohydrates are the primary lever after protein targets and fat floor are satisfied.
  \item Carbohydrates scale upward as training intensity and specificity increase (especially \textbf{RUN}, \textbf{RUCK}, and high-density \textbf{CAL} phases) and as test consequence increases.
  \item Carbohydrates scale upward during Deload/Test weeks by design (validity protection posture), using familiar, low GI-risk staples near protected windows.
  \item Novelty foods (new high-fiber loads, new sugar alcohols, new “performance” foods) are prohibited in the final 7–10 days before major gate tests.
  \end{itemize}

\newpage
\item Deload and test week nutrition posture aligned to Stage 3
  \begin{itemize}
  \item Default: raise calories up to maintenance during Deload/Test weeks to protect validity and recovery (this is a planned, repeatable posture, not a “diet break”).
  \item Maintain protein targets.
  \item Keep fat within the 20–30\% floor band.
  \item Preferentially raise carbohydrates to support performance and reduce fatigue contamination.
  \item No novelty foods and no untested supplement changes in the Deload/Test window.
  \end{itemize}

\item No novelty near gates
  \begin{itemize}
  \item No major diet structure change in the final 7–10 days before major gate tests (macro ratio shifts, new food classes, new timing schemes, new supplement categories).
  \item Any change \textbf{MUST} be logged with date, reason, and expected effect; changes are treated as controlled-variable events for validity interpretation.
  \end{itemize}

\item Budget and start-from-zero feasibility posture
  \begin{itemize}
  \item Default posture uses low-cost staples and minimal prep infrastructure; complexity is denied unless it improves adherence measurably.
  \item Stage 7.13 is the tooling layer for containers/hydration; Stage 10 does not prescribe meal plans or recipes.
  \end{itemize}
\end{itemize}

\newpage
\subsection*{10.A.2 Baseline Calculations and Starting Targets}
\phantomsection
\addcontentsline{toc}{subsection}{10.A.2 Baseline Calculations and Starting Targets}

This block is a transparent estimate using the fixed start-state; uncertainty is explicitly labeled.

\begin{enumerate}
\item Lean body mass estimate (uncertain)

\begin{itemize}
\item Inputs: bodyweight = 266 lb; body fat estimate = 40\% (uncertain).
\item LBM (lb) = bodyweight \ensuremath{\times} (1 - BF estimate)
  \begin{itemize}
  \item LBM $\approx$ 266 \ensuremath{\times} (1 - 0.40) = 266 \ensuremath{\times} 0.60 = 159.6 lb
  \end{itemize}
\item Uncertainty note: the BF estimate is approximate; LBM is therefore approximate and used only to set a conservative protein basis.
\end{itemize}

\item Maintenance calorie estimate by two methods (transparent, conservative)

Method 1: Mifflin–St Jeor male \textbf{BMR} \ensuremath{\times} conservative activity factor

\begin{itemize}
\item Inputs: age 35; height 175 cm; weight 121 kg
\item \textbf{BMR} = 10\ensuremath{\times}kg + 6.25\ensuremath{\times}cm - 5\ensuremath{\times}age + 5
  \begin{itemize}
  \item \textbf{BMR} = 10\ensuremath{\times}121 + 6.25\ensuremath{\times}175 - 5\ensuremath{\times}35 + 5
  \item \textbf{BMR} = 1210 + 1093.75 - 175 + 5
  \item \textbf{BMR} $\approx$ 2133.75 kcal/day ($\approx$ 2134)
  \end{itemize}
\item Conservative activity factor range (early pipeline with 6 sessions/week low-impact + steps baseline): 1.25–1.40
  \begin{itemize}
  \item Maintenance $\approx$ 2134 \ensuremath{\times} 1.25 to 2134 \ensuremath{\times} 1.40
  \item Maintenance $\approx$ 2668 to 2988 kcal/day
  \end{itemize}
\end{itemize}

Method 2: Bodyweight multiplier range (context-appropriate, coarse)

\begin{itemize}
\item Conservative range: 11–13 kcal/lb (recognizing obesity can inflate this estimate and low training age can deflate effective output)
  \begin{itemize}
  \item Maintenance $\approx$ 266 \ensuremath{\times} 11 to 266 \ensuremath{\times} 13
  \item Maintenance $\approx$ 2926 to 3458 kcal/day
  \end{itemize}
\end{itemize}

Working maintenance range (decision range)

\begin{itemize}
\item Overlap exists but is narrow: approximately 2926–2988 kcal/day.
\item Working maintenance range (widened conservatively due to method uncertainty and expected variability in steps/\textbf{NEAT}/season): 2900–3200 kcal/day.
  \begin{itemize}
  \item Rationale: method 1 may understate activity output if steps are consistently high; method 2 may overstate due to higher fat mass; early-phase variability is expected.
  \end{itemize}
\end{itemize}

\item Initial deficit range appropriate for Phase 00 (guardrail enforced)

\begin{itemize}
\item Weight-loss guardrail: planned average loss \textbf{MUST} NOT exceed 0.5\% bodyweight/week.
  \begin{itemize}
  \item 0.5\% of 266 lb = 1.33 lb/week (maximum planned average).
  \item Energy equivalent: 1.33 \ensuremath{\times} 3500 $\approx$ 4655 kcal/week \textrightarrow\ $\approx$ 665 kcal/day deficit ceiling (planned average).
  \end{itemize}
\item Phase 00 initial planned deficit range: 300–650 kcal/day below working maintenance.
  \begin{itemize}
  \item This is a planned range; weekly review using weekly bodyweight average governs whether it is held, reduced, or raised to maintain the guardrail.
  \end{itemize}
\newpage
\item Adherence and safety notes:
  \begin{itemize}
  \item If red-flag symptoms occur (dizziness/fainting, chest pain/pressure, syncope/near-syncope, severe GI distress, heat illness signs), stop and follow escalation posture; the plan is suspended until safe to resume.
  \item If weekly average loss exceeds the 0.5\% guardrail, reduce deficit immediately (move calories toward maintenance).
  \end{itemize}
\end{itemize}

\item Starting macro target example ranges (Phase 00 example; ranges, not prescriptions)

\begin{itemize}
\item Calories (Phase 00 start range): 2600–2900 kcal/day
  \begin{itemize}
  \item Derived from working maintenance 2900–3200 with an initial 300–650 kcal/day planned deficit, adjusted to remain conservative while monitoring weekly trends.
  \end{itemize}

\item Protein (g/day): hybrid anchor, 1.0–1.2 g/lb of chosen basis
  \begin{itemize}
  \item Anchor A basis (target \textbf{BW} 180): 180–216 g/day
  \item Anchor B basis (LBM $\approx$ 160): 160–192 g/day
  \item Hybrid rule selects the higher basis $\to$ Anchor A dominates early
  \item Working protein range (Phase 00): 180–210 g/day (upper bound held below the full 216 to preserve feasibility and reduce friction)
  \item Purpose: satiety, lean mass retention, recovery support.
  \end{itemize}

\item Fat floor: 20–30\% of calories (must remain within band)
  \begin{itemize}
  \item At 2600 kcal/day:
    \begin{itemize}
    \item 20\% fat = 520 kcal $\to$ 58 g/day (520 \ensuremath{\div} 9)
    \item 30\% fat = 780 kcal $\to$ 87 g/day (780 \ensuremath{\div} 9)
    \end{itemize}
  \item At 2900 kcal/day:
    \begin{itemize}
    \item 20\% fat = 580 kcal $\to$ 64 g/day
    \item 30\% fat = 870 kcal $\to$ 97 g/day
    \end{itemize}
  \item Working fat range (Phase 00): \textasciitilde 60–95 g/day, depending on where calories sit and whether appetite/adherence requires the higher end of the fat floor band.
  \end{itemize}

\item Carbs: remainder after protein and fat are set (primary lever)
  \begin{itemize}
  \item Calculation method:
    \begin{itemize}
    \item Carbs (g) = (Calories - (Protein\ensuremath{\times}4 + Fat\ensuremath{\times}9)) \ensuremath{\div} 4
    \end{itemize}
  \item Example remainder bands (Phase 00, using typical midpoints):
    \begin{itemize}
    \item If 2750 kcal, protein 190 g (760 kcal), fat 75 g (675 kcal): carbs $\approx$ (2750 - 1435) \ensuremath{\div} 4 $\approx$ 329 g/day
    \item If 2600 kcal, protein 200 g (800 kcal), fat 85 g (765 kcal): carbs $\approx$ (2600 - 1565) \ensuremath{\div} 4 $\approx$ 259 g/day
    \end{itemize}
  \item Working carb range (Phase 00): \textasciitilde 220–340 g/day, adjusted primarily by training tolerance, GI stability, and adherence, while maintaining “no novelty near gates” posture.
  \end{itemize}
\end{itemize}

\end{enumerate}

\begin{landscape}

\subsection*{10.A.3 Phase-by-Phase Macro Ladder Table}
\phantomsection
\addcontentsline{toc}{subsection}{10.A.3 Phase-by-Phase Macro Ladder Table}

\begingroup
\scriptsize
\setlength{\tabcolsep}{2pt}
\renewcommand{\arraystretch}{1.12}

\begin{longtblr}{
  width=\linewidth,
  rowhead=1,
  colspec = {
    X[l,2.10]
    Q[l,wd=0.35in]
    X[l,3.50]
    X[l,2.45]
    X[l,2.35]
    Q[l,wd=0.5in]
    Q[l,wd=0.5in]
    X[l,2.60]
    X[l,2.60]
    X[l,2.85]
  },
  hlines,
  vlines,
  cells = {hsep=6pt, vsep=6pt, halign=l, valign=t},
  row{odd} = {bg=white},
  row{even} = {bg=LightGray},
  row{1} = {bg=StageBlue!15, fg=black, font=\bfseries, halign=l, valign=t},
}
\Hdr{Phase} &
\Hdr{Duration\\(Weeks)} &
\Hdr{Nutrition \\Mission} &
\Hdr{Bodyweight \\Trend Target} &
\Hdr{Calorie \\Posture} &
\Hdr{Protein \\Target} &
\Hdr{Fat Floor \\Rule} &
\Hdr{Carb \\Strategy} &
\Hdr{Deload/Test Week \\Modification} &
\Hdr{Notes} \\
Phase 00 — Bodily Reactivation and Baseline Creation &
16 &
Establish deficit tolerance, tracking discipline, and safe weight-loss trend without performance chasing &
0.25–0.5\% \textbf{BW}/week average; never >0.5\% &
Moderate deficit (typically maintenance -300 to -650 kcal/day) &
180–210 g/day &
20–30\% of calories &
Moderate carbs as remainder; prefer stable, familiar staples; GI stability prioritized &
Raise to maintenance; protein stable; fat within band; carbs upshift; no novelty &
Start-from-zero logistics; 7.13 minimal prep posture. Avoid aggressive deficit due to low training age and higher risk. \\
Phase 01 — Base Conditioning and Hypertrophy &
24 &
Continue body comp change while protecting recovery as training volume increases &
0.25–0.5\% \textbf{BW}/week average; never >0.5\% &
Moderate deficit (typically -250 to -500) &
175–205 g/day &
20–30\% &
Carbs slightly higher than Phase 00 to support volume; stable sources; no novelty near Deload/Test &
Raise to maintenance; carb upshift; avoid new foods and new supplement categories &
Avoid punitive restriction; consistency drives adherence. Supplements optional category references only (7.03–7.06). \\
Phase 02 — Maximal Strength Development &
12 &
Protect strength development and recovery; slow fat loss if needed &
0.15–0.35\% \textbf{BW}/week average; never >0.5\% &
Mild deficit to near maintenance (typically -0 to -300) &
170–205 g/day &
20–30\% &
Carbs increase to support strength sessions; keep fats within floor; maintain stable timing &
Raise to maintenance; carb upshift preferentially; no novelty &
Aggressive deficit avoided to protect strength and tissue tolerance; maintain tracking discipline. \\
Phase 03 — Converting Strength to Power and Endurance for \textbf{SOF} Selection &
12 &
Maintain performance while continuing controlled fat loss &
0.15–0.35\% \textbf{BW}/week average; never >0.5\% &
Mild deficit to near maintenance (-0 to -250) &
170–200 g/day &
20–30\% &
Carbs support mixed demands; stabilize pre-gate foods 7–10 days prior &
Raise to maintenance; carbs upshift; no novelty &
Phase begins to increase consequence; preserve test validity by stabilizing diet variables. \\
Phase 04 — Strength-Endurance and Work Capacity Development &
12 &
Support work capacity and calisthenics volume without recovery collapse &
0.15–0.35\% \textbf{BW}/week average; never >0.5\% &
Mild deficit (-150 to -350) &
170–200 g/day &
20–30\% &
Carbs used to prevent performance degradation; stable staples; avoid new high-fiber loads &
Raise to maintenance; carbs upshift; no novelty &
If \textbf{SWIM} becomes active, avoid GI-risk changes near swim test windows; preserve validity. \\
Phase 05 — Aerobic Base and Extended Stamina Development &
12 &
Protect endurance development; avoid low-energy availability &
0.15–0.30\% \textbf{BW}/week average; never >0.5\% &
Mild deficit to near maintenance (-100 to -300) &
165–200 g/day &
20–30\% &
Carbs increase to support longer sessions; keep familiar low GI-risk carbs near protected windows &
Raise to maintenance; carbs upshift; no novelty &
Avoid aggressive deficit because endurance and recovery demands rise; maintain stable hydration practices (no extremes). \\
Phase 06 — Lactate Threshold and Power-Endurance Development &
12 &
Prioritize performance and recovery; fat loss becomes secondary &
0.0–0.25\% \textbf{BW}/week average; never >0.5\% &
Near maintenance (-0 to -200) &
165–195 g/day &
20–30\% &
Carbs higher; support higher intensity; stabilize food choices near gates &
Raise to maintenance; carb upshift; no novelty &
This phase often exposes validity drift from under-fueling; hold conservative posture to protect gate evidence. \\
Phase 07 — Advanced Hybrid Overload &
12 &
Sustain high density without over-restriction; preserve recovery &
0.0–0.25\% \textbf{BW}/week average; never >0.5\% &
Near maintenance (-0 to -150) &
165–195 g/day &
20–30\% &
Carbs remain high and stable; avoid large swings; no novelty near protected windows &
Raise to maintenance; carb upshift; no novelty &
Over-restriction risk is high; treat consistency as the control. \\
Phase 08 — Structural Durability and Load Tolerance Phase &
12 &
Protect durability under load; avoid diet-driven injury risk &
0.0–0.25\% \textbf{BW}/week average; never >0.5\% &
Near maintenance (-0 to -150) &
165–195 g/day &
20–30\% &
Carbs support load tolerance; stabilize pre-ruck windows; avoid GI novelty &
Raise to maintenance; carb upshift; no novelty &
\textbf{RUCK} tolerance is high consequence; under-fueling increases injury risk and invalid evidence risk. \\
Phase 09 — High-Volume Calisthenics and Stamina &
12 &
Support high \textbf{CAL} volume; protect soft tissue and recovery &
0.0–0.25\% \textbf{BW}/week average; never >0.5\% &
Near maintenance (-0 to -150) &
160–195 g/day &
20–30\% &
Carbs support \textbf{CAL} density; stable routine; avoid new stimulant/supplement changes near gates &
Raise to maintenance; carb upshift; no novelty &
\textbf{CAL} volume is sensitive to low energy availability; preserve sleep and recovery markers. \\
Phase 10 — Advanced Tactical Readiness Physical &
12 &
Maintain performance and test readiness; minimize variance &
0.0–0.20\% \textbf{BW}/week average; never >0.5\% &
Maintenance to slight surplus (-0 to +150) &
160–190 g/day &
20–30\% &
Carbs high and stable; no novelty near gates; emphasize repeatability &
Raise to maintenance (or hold if already at maintenance); carb upshift; no novelty &
Underwater standards are governance-gated; do not change diet variables near any high-consequence pool windows. \\
Phase 11 — Pre-Selection Conditioning and Recovery &
12 &
Protect recovery; maintain readiness; avoid diet fatigue &
0.0–0.20\% \textbf{BW}/week average; never >0.5\% &
Maintenance to slight deficit (-150 to -0) &
160–190 g/day &
20–30\% &
Carbs stable; prioritize GI predictability; avoid new foods &
Raise to maintenance; carb upshift; no novelty &
Recovery emphasis increases; under-fueling increases test invalidation risk via fatigue contamination. \\
Phase 12 — Final Pre-Selection Conditioning &
12 &
Maximize repeatable output; minimize confounders &
0.0–0.15\% \textbf{BW}/week average; never >0.5\% &
Maintenance to slight surplus (0 to +200) &
160–190 g/day &
20–30\% &
Carbs high; strict no novelty; pre-window stabilization is mandatory &
Maintain or slight upshift as needed; carb upshift; no novelty &
Any major diet structure change is prohibited in the final 7–10 days before major gate tests; planned deload posture is allowed but must remain consistent. \\
Phase 13 — Full-Spectrum Readiness Consolidation &
12 &
Consolidate selection-ready envelope with low variance &
0.0–0.15\% \textbf{BW}/week average; never >0.5\% &
Maintenance to slight surplus (0 to +200) &
160–190 g/day &
20–30\% &
Carbs high and stable; focus on repeatability; no novelty &
Maintain or slight upshift as needed; carb upshift; no novelty &
End-state verification depends on stable nutrition variables; avoid last-minute restriction. \\
\end{longtblr}

\endgroup

\end{landscape}

\subsection*{10.A.4 Minimal vs Ideal Nutrition Execution Posture}
\phantomsection
\addcontentsline{toc}{subsection}{10.A.4 Minimal vs Ideal Nutrition Execution Posture}

Minimal Path (budget constrained)

\begin{itemize}
\item Feasibility posture:
  \begin{itemize}
  \item rely on simple staples with limited variety to reduce decision fatigue and cost,
  \item minimal prep infrastructure consistent with 7.13 Tier 1 posture (containers, basic hydration container, cleaning tools).
  \end{itemize}
\item Friction-reduction posture:
  \begin{itemize}
  \item consistency over complexity,
  \item repeatable meal timing and repeatable food selection,
  \item one change at a time when adjustments are required (logged with date/reason/expected effect).
  \end{itemize}
\item Optional supplement category references only (no dosing):
  \begin{itemize}
  \item 7.03 / 7.04 / 7.05 / 7.06 are optional supports, never substitutes for sleep, training dose control, or food adequacy.
  \end{itemize}
\end{itemize}

Ideal Path (more resources)

\begin{itemize}
\item What improves adherence and quality without changing macro rules:
  \begin{itemize}
  \item greater food variety while preserving GI predictability,
  \item higher convenience options that reduce prep time and missed meals,
  \item expanded 7.13 tooling (more containers, better carry insulation, improved workflow consistency).
  \end{itemize}
\item Still governed by:
  \begin{itemize}
  \item the 0.5\% \textbf{BW}/week guardrail,
  \item no planned refeeds/diet breaks,
  \item no novelty near gates/test weeks,
  \item one change at a time posture,
  \item logging discipline and weekly review discipline.
  \end{itemize}
\end{itemize}

\newpage
\subsection*{10.A.5 Macro Ladder Integrity Checks}
\phantomsection
\addcontentsline{toc}{subsection}{10.A.5 Macro Ladder Integrity Checks}

\begin{itemize}
\item Early phases prioritize safe deficit and habit formation within the 0.5\% \textbf{BW}/week guardrail; planned average loss never exceeds 0.5\% \textbf{BW}/week.
\item Mid phases explicitly moderate deficit to protect \textbf{STR} development, \textbf{RUN} tolerance, and recovery; aggressive restriction is denied when it would degrade performance evidence.
\item Late phases default to near-maintenance to slight surplus posture with guardrails to protect specificity and test validity; weight loss becomes secondary to repeatable output.
\item Deload/Test weeks raise calories to maintenance by design, with protein stable and fat within the floor band, and carbs upshifted; this posture is planned and repeatable, not “novelty.”
\item No major diet structure changes are introduced in the final 7–10 days before major gates; any change is logged with date, reason, expected effect.
\item Targets are measurable using Stage 2 logs:
  \begin{itemize}
  \item weekly bodyweight average (\textbf{MET-2B-001}),
  \item weekly circumference suite (\textbf{MET-2B-003} through \textbf{MET-2B-009}),
  \item recovery markers (\textbf{MET-2B-038} through \textbf{MET-2B-042}).
  \end{itemize}
\item Feasibility matches start-from-zero posture and 7.13 tooling assumptions; no meal plan dependence and no specialty procurement.
\item Supplement references remain category-level optional only (7.03–7.06); 7.02 remains education-only and never required.
\item Red-flag symptoms override the plan; the nutrition posture is suspended and escalation posture is followed when safety signals appear.
\item Any detected conflict with locked upstream doctrine is flagged for Stage 8 review; no silent edits are performed.
\end{itemize}

\end{document}
